\subsection{The canonical free resolution}

For every A-module M, denote by \(\mathbf{L}_{0}(\mathbf{M})\) the free A-module \(\mathbf{A}^{(\mathbf{M})}\) with basis M, \((e_{m})_{m\in\mathbf{M}}\) its canonical basis and \(p_{\mathbf{M}}:\mathbf{L}_{0}(\mathbf{M})\to\mathbf{M}\) the homomorphism such that

\[
p _ {\mathbf {M}} (e _ {m}) = m, \quad m \in \mathbf {M}.
\]

Set \(Z_0(\mathbf{M}) = \operatorname{Ker} p_{\mathbf{M}}\) and let \(i_{\mathbf{M}}: Z_0(\mathbf{M}) \to L_0(\mathbf{M})\) be the canonical injection. We have an exact sequence

\[
0 \longrightarrow Z _ {0} (\mathrm{M}) \stackrel {i _ {\mathrm{M}}} {\longrightarrow} L _ {0} (\mathrm{M}) \stackrel {p _ {\mathrm{M}}} {\longrightarrow} \mathrm{M} \longrightarrow 0.\tag{1}
\]

We define a graded module L(M) by setting \(\mathrm{L}_{n}(\mathrm{M}) = 0\) for n \textless{} 0 and, by induction on the integer n \textgreater{} 0

\[
\mathrm{L} _ {n} (\mathrm{M}) = \mathrm{L} _ {0} \left(\mathrm{Z} _ {n - 1} (\mathrm{M})\right); \quad \mathrm{Z} _ {n} (\mathrm{M}) = \mathrm{Z} _ {0} \left(\mathrm{Z} _ {n - 1} (\mathrm{M})\right).\tag{2}
\]

We define A-homomorphisms \(d_{n}^{\mathbf{M}}:\mathrm{L}_{n}(\mathbf{M})\to \mathrm{L}_{n - 1}(\mathbf{M})\) by

\[
\left\{ \begin{array}{l l} d _ {n} ^ {\mathbf {M}} = 0, & n \leqslant 0, \\ d _ {1} ^ {\mathbf {M}} = i _ {\mathbf {M}} \circ p _ {\mathbf {Z} _ {0} (\mathbf {M})}, & \\ d _ {n} ^ {\mathbf {M}} = i _ {\mathbf {Z} _ {n - 2} (\mathbf {M})} \circ p _ {\mathbf {Z} _ {n - 1} (\mathbf {M})}, & n > 1. \end{array} \right.\tag{3}
\]

By construction, we have an exact sequence

\[
\longrightarrow \mathrm{L} _ {n} (\mathrm{M}) \xrightarrow {d _ {n} ^ {\mathrm{M}}} \mathrm{L} _ {n - 1} (\mathrm{M}) \longrightarrow \dots \longrightarrow \mathrm{L} _ {0} (\mathrm{M}) \xrightarrow {p _ {\mathrm{M}}} \mathrm{M} \longrightarrow 0,
\]

so that, if one extends \(p_{\mathbf{M}}\) as a morphism of complexes

\[
p _ {\mathbf {M}}: (\mathrm{L} (\mathbf {M}), d ^ {\mathbf {M}}) \rightarrow \mathbf {M},
\]

one obtains a free resolution of M, called the canonical free resolution of M.

Let \(f: \mathbf{M} \to \mathbf{N}\) be a homomorphism of A-modules. Denote by

\[
\mathrm{L} _ {0} (f): \mathrm{L} _ {0} (\mathrm{M}) \rightarrow \mathrm{L} _ {0} (\mathrm{N})
\]

the unique A-homomorphism such that \(L_0(f)(e_m) = e_{f(m)}\) for all \(m \in \mathbf{M}\) . We have

\[
p _ {\mathbf {N}} \circ \mathrm{L} _ {0} (f) = f \circ p _ {\mathbf {M}}.\tag{4}
\]

Consequently, \(L_0(f)\) induces an A-homomorphism \(Z_0(f):Z_0(M)\to Z_0(N)\) and we have

\[
i _ {\mathbf {N}} \circ \mathbf {Z} _ {0} (f) = \mathbf {L} _ {0} (f) \circ i _ {\mathbf {M}}.\tag{5}
\]

Set \(\mathrm{L}_{n}(f)=0\) , for n\textless0 and define by induction on the integer n\textgreater0, homomorphisms \(\mathrm{L}_{n}(f):\mathrm{L}_{n}(\mathrm{M})\to\mathrm{L}_{n}(\mathrm{N})\) and \(\mathrm{Z}_{n}(f):\mathrm{Z}_{n}(\mathrm{M})\to\mathrm{Z}_{n}(\mathrm{N})\) by

\[
\left\{ \begin{array}{l} \mathrm{L} _ {n} (f) = \mathrm{L} _ {0} (\mathrm{Z} _ {n - 1} (f)) \\ \mathrm{Z} _ {n} (f) = \mathrm{Z} _ {0} (\mathrm{Z} _ {n - 1} (f)). \end{array} \right.\tag{6}
\]

PROPOSITION 4. --- L(f) : L(M) → L(N) is a morphism of complexes of A-modules; moreover, we have \(p_{\mathbf{M}} \circ \mathrm{L}(f) = f \circ p_{\mathbf{N}}\) .

We need to prove that, for every integer n \textgreater{} 0 , the formula

\[
d _ {n} ^ {\mathbf {N}} \circ \mathrm{L} _ {n} (f) = \mathrm{L} _ {n - 1} (f) \circ d _ {n} ^ {\mathbf {M}}.
\]

We first have

\[
\begin{array}{r l r l} d _ {1} ^ {\mathbf {N}} \circ \mathrm{L} _ {1} (f) & = i _ {\mathrm{N}} \circ p _ {\mathbf {Z} _ {0} (\mathbf {N})} \circ \mathrm{L} _ {0} (\mathbf {Z} _ {0} (f)) & & (\mathrm {by~(3)~and~(6)}) \\ & = i _ {\mathrm{N}} \circ \mathrm{Z} _ {0} (f) \circ p _ {\mathbf {Z} _ {0} (\mathbf {M})} & & (\mathrm {by~(4)}) \\ & = \mathrm{L} _ {0} (f) \circ i _ {\mathbf {M}} \circ p _ {\mathbf {Z} _ {0} (\mathbf {M})} & & (\mathrm {by~(5)}) \\ & = \mathrm{L} _ {0} (f) \circ d _ {1} ^ {\mathbf {M}} & & (\mathrm {by~(3))}. \end{array}
\]

When \(n > 1\) , we successively have

\[
\begin{array}{l l} d _ {n} ^ {\mathbf {N}} \circ \mathrm{L} _ {n} (f) = i _ {\mathrm{Z} _ {n - 2 (\mathrm{N})}} \circ p _ {\mathrm{Z} _ {n - 1 (\mathrm{N})}} \circ \mathrm{L} _ {0} (\mathrm{Z} _ {n - 1} (f)) & \text {(by (3) and (6))} \\ = i _ {\mathrm{Z} _ {n - 2 (\mathrm{N})}} \circ \mathrm{Z} _ {n - 1} (f) \circ p _ {\mathrm{Z} _ {n - 1} (\mathrm{M})} & \text {(by (4))} \\ = i _ {\mathrm{Z} _ {n - 2 (\mathrm{N})}} \circ \mathrm{Z} _ {0} (\mathrm{Z} _ {n - 2} (f)) \circ p _ {\mathrm{Z} _ {n - 1} (\mathrm{M})} & \text {(by (6))} \\ = \mathrm{L} _ {0} (\mathrm{Z} _ {n - 2} (f)) \circ i _ {\mathrm{Z} _ {n - 2 (\mathrm{M})}} \circ p _ {\mathrm{Z} _ {n - 1} (\mathrm{M})} & \text {(by (5))} \\ = \mathrm{L} _ {n - 1} (f) \circ d _ {n} ^ {\mathbf {M}} & \text {(by (3) and (6)) .} \end{array}
\]

We immediately have

\[
\mathrm{L} (1 _ {\mathbf {M}}) = 1 _ {\mathrm{L} (\mathbf {M})}.\tag{7}
\]

On the other hand, if \(g: \mathbf{N} \to \mathbf{P}\) is a homomorphism of A-modules, we have

\[
\mathrm{L} (g \circ f) = \mathrm{L} (g) \circ \mathrm{L} (f).\tag{8}
\]

Indeed, we have for \(m \in \mathbf{M}\) ,

\[
\mathrm{L} _ {0} (g \circ f) \left(e _ {m}\right) = e _ {g \circ f (m)} = \mathrm{L} _ {0} (g) \left(e _ {f (m)}\right) = \mathrm{L} _ {0} (g) \circ \mathrm{L} _ {0} (f) \left(e _ {m}\right),
\]

hence \(\mathrm{L}_{0}(g\circ f)=\mathrm{L}(g)\circ\mathrm{L}(f)\) ; therefore \(\mathrm{Z}_{0}(g\circ f)=\mathrm{Z}_{0}(g)\circ\mathrm{Z}_{0}(f)\) ; whence immediately \(\mathrm{L}_{n}(g\circ f)=\mathrm{L}_{n}(g)\circ\mathrm{L}_{n}(f)\) , for \(n\geqslant0\) , by induction on n, whence (8).

Remark. --- If f, g ∈ Hom\_A (M, N), one does not have L(f + g) = L(f) + L(g). However, these two morphisms are homotopic according to X, p.~49, prop. 3.

Let M be a right A-module; denote by A° the opposite ring of A, by M° the underlying A°-module of M, and by L(M°) its canonical free resolution. We denote by L(M) and call the canonical free resolution of M the underlying A-complex L(M°)° of L(M°). We thus have

\[
\mathrm{L} (\mathbf {M} ^ {\circ}) = \mathrm{L} (\mathbf {M}) ^ {\circ}.
\]
% semantic-fragments: legacy-input-seam
\relax{}
